\section{Theory Task T4: Alignment-Defect Decomposition and Quantum Schedules}
\label{sec:theory-t4-defect-schedule}

T3 provides a conditional convergence interface and T3a shows that aggregate
alignment is usually positive on the frozen MLP.  It does not identify which
part of the client--encoder--server chain produces the exceptional alignment
defect, nor whether replacing the empirical quantum schedule by an
asymptotically admissible one improves the finite run.  T4 addresses these two
questions without changing the frozen baseline claim.

\begin{takeawaybox}[title={T4 result in one sentence}]
The defect admits an exact client-pulse decomposition, but its adverse
heterogeneity term is cancelled by comparably large favorable local-update
terms, making a termwise absolute bound unusably loose; an admissible
$r^{-0.6}$ quantum does not repair late-stage descent and either reduces
accuracy or, after horizon matching, reduces worst-class accuracy.
\end{takeawaybox}

\subsection{Exact client-pulse decomposition}

Let $\mathcal{S}_i^r$ be the coordinates emitted by client $i$ in round $r$,
and let $s_{ij}^r\in\{-1,+1\}$ denote an emitted sign.  Write
\begin{equation}
    g_i^r=\nabla F_i(w^r),
    \qquad
    g^r=\sum_{i=1}^Mp_i g_i^r,
    \qquad
    v_i^r=\frac{\delta_i^r}{\eta E}.
    \label{eq:t4-local-proxy}
\end{equation}
The proxy $v_i^r$ has the sign of the current model delta and approximates
$-g_i^r$ when local drift and stochastic-gradient error are small.  Define
\begin{align}
    P_r
    &=\sum_i\sum_{j\in\mathcal{S}_i^r}|v_{ij}^r|,
    \label{eq:t4-ideal-local-mass}\\
    R_r
    &=2\sum_i\sum_{j\in\mathcal{S}_i^r}|v_{ij}^r|
      \mathbf{1}_{\{s_{ij}^rv_{ij}^r<0\}},
    \label{eq:t4-memory-penalty}\\
    L_r
    &=\sum_i\sum_{j\in\mathcal{S}_i^r}
      (-g_{ij}^r-v_{ij}^r)s_{ij}^r,
    \label{eq:t4-drift-term}\\
    B_r
    &=\sum_i\sum_{j\in\mathcal{S}_i^r}
      (g_{ij}^r-g_j^r)s_{ij}^r.
    \label{eq:t4-heterogeneity-term}
\end{align}
Here $P_r$ is the ideal alignment with the current local-update proxy,
$R_r$ is the penalty from accumulated state emitting against that current
proxy, $L_r$ combines local-SGD drift and stochastic-gradient mismatch, and
$B_r$ is the client-to-global heterogeneity term.  The signs of $L_r$ and
$B_r$ are not fixed.

\begin{proposition}[Exact T4 alignment decomposition]
\label{prop:t4-exact-decomposition}
For the aggregate pulse $C^r=\sum_i c_i^r$, the T3 net alignment satisfies
\begin{equation}
    A_r^{\mathrm{net}}
    =-\langle g^r,C^r\rangle
    =P_r-R_r+L_r+B_r.
    \label{eq:t4-exact-decomposition}
\end{equation}
\end{proposition}

\begin{proof}
On every emitted client-coordinate pair,
\begin{equation}
    -g_j^rs_{ij}^r
    =v_{ij}^rs_{ij}^r
     +(-g_{ij}^r-v_{ij}^r)s_{ij}^r
     +(g_{ij}^r-g_j^r)s_{ij}^r.
\end{equation}
Moreover,
$v_{ij}^rs_{ij}^r=|v_{ij}^r|$ when the signs agree and
$v_{ij}^rs_{ij}^r=-|v_{ij}^r|$ when they disagree.  Summing gives
\eqref{eq:t4-exact-decomposition}.
\end{proof}

For a fixed $\kappa>0$, define the coverage deficit
\begin{equation}
    Q_r(\kappa)
    =\left[\kappa\|g^r\|_2^2-P_r\right]_+.
    \label{eq:t4-coverage-deficit}
\end{equation}
It captures the part of the target that cannot be supplied by the emitted
current-update mass.  The non-emitted proxy mass
\begin{equation}
    V_r^{\mathrm{silent}}
    =\sum_i\sum_{j\notin\mathcal{S}_i^r}|v_{ij}^r|
\end{equation}
is reported separately as a deadzone/silence diagnostic.

\begin{corollary}[Nonnegative defect upper bound]
\label{cor:t4-defect-upper-bound}
The realized T3 defect obeys
\begin{equation}
    \beta_r(\kappa)
    \leq
    Q_r(\kappa)+R_r+[-L_r]_++[-B_r]_+.
    \label{eq:t4-defect-upper-bound}
\end{equation}
\end{corollary}

\begin{proof}
Insert \eqref{eq:t4-exact-decomposition} into
\eqref{eq:t3-audit-defect} and use
$[a+b+c+d]_+\leq[a]_++[b]_++[c]_++[d]_+$, noting that $R_r\geq0$.
\end{proof}

The bound is deliberately one-sided.  Its components need not be a causal or
additive attribution of the observed $\beta_r$, because large signed terms in
\eqref{eq:t4-exact-decomposition} may cancel before the positive part is
taken.  T4 measures both the exact signed identity and the conservative bound.

\subsection{Schedule experiment and observer-free protocol}

The experiment uses the same strong non-IID MLP, three held-out partition
seeds, training seeds, and 31 audit rounds as T3a.  Unlike T3a's
$5000$-example global diagnostic, T4 evaluates every $g_i^r$ on all $6000$
examples assigned to client $i$.  Consequently, $g^r$ is the exact gradient of
the empirical federated training objective for that partition.

The trigger is numerically hybrid: one last-bit change near $\vartheta$ can
change a pulse and then the later trajectory.  In a pilot sequential audit,
the additional reference-gradient work changed one seed's first pulse count at
round $40$.  The final protocol therefore audits each target round by an
independent replay from initialization.  The pulse decision is completed and
copied before any reference-gradient evaluation.  Final performance is taken
from a separate uninstrumented run.  This removes the diagnostic observer
effect up to every audited pulse.

Three quantum schedules are compared:
\begin{center}
\begin{tabular}{lcccc}
\toprule
Schedule & $q_0$ & $\alpha$ & $\sum_{r=1}^{150}q_r$ & Asymptotic window\\
\midrule
Frozen baseline & $0.005000$ & $0.1$ & $0.711509$ & no\\
Admissible, raw & $0.005000$ & $0.6$ & $0.552320$ & yes\\
Admissible, matched & $0.006441$ & $0.6$ & $0.711509$ & yes\\
\bottomrule
\end{tabular}
\end{center}
The raw schedule tests the direct replacement $\alpha=0.1\mapsto0.6$.  The
matched schedule rescales $q_0$ to preserve the finite-horizon quantum sum, so
that a change cannot be attributed merely to a $22.4\%$ smaller update budget.

\subsection{T4 results}

\begin{table}[H]
\centering
\caption{T4 schedule comparison over three held-out seeds.  Values are
mean $\pm$ sample standard deviation.  ``Late'' denotes rounds $101$--$150$;
alignment and descent fractions use the 31 independent audit snapshots.}
\label{tab:t4-schedule-summary}
\resizebox{\textwidth}{!}{%
\begin{tabular}{lrrrrrrrr}
\toprule
Schedule & Test acc. & Worst class & $\widehat\kappa_R$ & Late align. & $\Delta F<0$ & Late $\Delta F<0$ & Events\\
\midrule
Baseline
& $83.01\!\pm\!0.19\%$ & $59.30\!\pm\!3.61\%$
& $35.42\!\pm\!1.35$ & $27.66\!\pm\!3.07$
& $75.27\!\pm\!7.45\%$ & $73.33\!\pm\!15.28\%$
& $0.990\!\pm\!0.011$ M\\
Admissible, raw
& $82.40\!\pm\!0.12\%$ & $54.37\!\pm\!2.74\%$
& $38.19\!\pm\!0.61$ & $27.67\!\pm\!2.72$
& $79.57\!\pm\!6.72\%$ & $63.33\!\pm\!15.28\%$
& $1.004\!\pm\!0.013$ M\\
Admissible, matched
& $82.94\!\pm\!0.10\%$ & $54.13\!\pm\!0.76\%$
& $36.24\!\pm\!3.10$ & $29.05\!\pm\!4.43$
& $68.82\!\pm\!7.45\%$ & $63.33\!\pm\!15.28\%$
& $0.971\!\pm\!0.019$ M\\
\bottomrule
\end{tabular}}
\end{table}

\begin{figure}[H]
\centering
\begin{tikzpicture}
\begin{groupplot}[
  group style={group size=2 by 1, horizontal sep=1.35cm},
  width=0.46\linewidth,
  height=0.31\linewidth,
  xlabel={Round},
  grid=major,
  major grid style={black!12},
  tick label style={font=\small},
  label style={font=\small},
  title style={font=\small},
]
\nextgroupplot[
  title={Mean aggregate alignment},
  ylabel={$A_r^{\mathrm{net}}/\|g^r\|_2^2$},
  legend entries={baseline,admissible,matched},
  legend style={font=\scriptsize, draw=none, at={(0.98,0.98)}, anchor=north east},
]
\addplot+[mark=*, mark size=1.1pt] table[x=round,y=alignment_ratio,col sep=comma]
  {data/t4_schedule_baseline.csv};
\addplot+[mark=square*, mark size=1.0pt] table[x=round,y=alignment_ratio,col sep=comma]
  {data/t4_schedule_admissible.csv};
\addplot+[mark=triangle*, mark size=1.1pt] table[x=round,y=alignment_ratio,col sep=comma]
  {data/t4_schedule_admissible_matched.csv};
\addplot[black, thin, forget plot] coordinates {(1,0) (150,0)};

\nextgroupplot[
  title={Late objective increments},
  ylabel={$F(w^{r+1})-F(w^r)$},
  xmin=50,
  ymin=-0.013,
  ymax=0.013,
]
\addplot+[mark=*, mark size=1.1pt] table[x=round,y=objective_change,col sep=comma]
  {data/t4_schedule_baseline.csv};
\addplot+[mark=square*, mark size=1.0pt] table[x=round,y=objective_change,col sep=comma]
  {data/t4_schedule_admissible.csv};
\addplot+[mark=triangle*, mark size=1.1pt] table[x=round,y=objective_change,col sep=comma]
  {data/t4_schedule_admissible_matched.csv};
\addplot[black, thin, forget plot] coordinates {(50,0) (150,0)};
\end{groupplot}
\end{tikzpicture}
\caption{T4 schedule comparison.  Faster decay raises the full-horizon
alignment statistic but does not improve late alignment or make late objective
increments reliably negative.}
\label{fig:t4-schedule-comparison}
\end{figure}

Table~\ref{tab:t4-schedule-summary} and
Figure~\ref{fig:t4-schedule-comparison} show that asymptotic admissibility is
not a finite-horizon remedy by itself.
\begin{enumerate}
    \item The raw $\alpha=0.6$ schedule reduces the quantum sum by $22.4\%$.
    It loses $0.61$ percentage points of mean test accuracy, raises the final
    train objective by $0.0182$, and does not change the mean late alignment
    ratio ($27.67$ versus $27.66$).  Its late descent fraction is lower by ten
    percentage points.
    \item Matching the quantum sum restores mean test accuracy to within
    $0.07$ percentage points and reduces events by $1.9\%$, but it still lowers
    the late descent fraction by ten percentage points.  More importantly, its
    worst-class accuracy is lower on all three seeds and falls from $59.30\%$
    to $54.13\%$ on average.
    \item The small differences in mean test accuracy between the baseline and
    matched schedules are below the numerical branch variability exposed by
    the hybrid trigger.  The consistent worst-class loss and the absence of a
    late-descent gain are the more reliable conclusions.
\end{enumerate}

\subsection{What produces the measured defect?}

For comparability with T3a, set $\kappa=8$.  The following quantities pool all
three seeds and normalize the $q_r$-weighted sums by
$\sum_rq_r\kappa\|g^r\|_2^2$.

\begin{table}[H]
\centering
\caption{Exact signed T4 decomposition and conservative defect bound.  Large
positive and negative entries demonstrate cancellation; the final two rows
must not be read as additive attribution shares.}
\label{tab:t4-defect-decomposition}
\begin{tabular}{lrrr}
\toprule
Normalized quantity & Baseline & Admissible, raw & Admissible, matched\\
\midrule
Ideal local mass $P_r$ & $39.002$ & $43.058$ & $36.932$\\
Memory penalty $R_r$ & $0.000$ & $0.000$ & $0.000$\\
Local drift/noise $L_r$ & $+49.864$ & $+54.506$ & $+46.717$\\
Heterogeneity $B_r$ & $-84.445$ & $-92.789$ & $-79.132$\\
Net alignment $A_r^{\mathrm{net}}$ & $4.420$ & $4.775$ & $4.517$\\
\midrule
Realized defect $\beta_r(8)$ & $0.126$ & $0.160$ & $0.144$\\
Upper bound in \eqref{eq:t4-defect-upper-bound} & $84.738$ & $93.190$ & $79.480$\\
\bottomrule
\end{tabular}
\end{table}

The exact identity holds to a maximum error below $5\times10^{-14}$.  On the
baseline, every one of the $93$ snapshots has $R_r=0$: whenever a coordinate
fires, its state sign agrees with the current local model-delta sign.  Likewise,
$Q_r(8)=0$ on all $93$ snapshots.  Only about $7.58\%$ of the current local
proxy $\ell_1$ mass lies on emitted coordinates, but that selected mass still
exceeds $8\|g^r\|_2^2$ in every snapshot.  Thus neither memory-opposition nor
insufficient emitted proxy mass explains the observed $\kappa=8$ violations in
this regime.

The adverse signed term is client-to-global heterogeneity.  However, it is
almost cancelled by the ideal local mass and the positive realized
drift/noise term.  Dropping this cancellation produces an upper bound roughly
$670$ times larger than the realized baseline defect ($84.738/0.126$).  A proof
that bounds heterogeneity, local drift, and useful local alignment separately
by absolute values will therefore be far too conservative.  The next analysis
must preserve their covariance or directly control the combined quantity
$P_r+L_r+B_r$.

\subsection{T4 conclusions and next theorem target}

T4 closes the requested diagnostic step but not the convergence theorem:
\begin{enumerate}
    \item the aggregate alignment defect has an exact, verified decomposition
    into current local-update mass, memory opposition, local drift/noise, and
    heterogeneity;
    \item for the frozen strong non-IID MLP, stale-state sign opposition and
    the $\kappa=8$ coverage deficit are absent at the audited pulses;
    \item heterogeneity is the dominant adverse signed contribution, but large
    cancellation makes the nonnegative componentwise bound vacuous;
    \item placing $\alpha$ in the Robbins--Monro window is mathematically
    necessary for the T3 asymptotic corollary but does not improve late descent
    in the tested finite horizon;
    \item horizon matching avoids the raw schedule's overall-accuracy loss but
    exposes a consistent worst-class penalty.
\end{enumerate}

The next theorem target should therefore be a coupled alignment bound rather
than four independent worst-case bounds.  The controlled $2\times2$
decomposition over IID versus strong non-IID partitions and $E=1$ versus
$E=5$ has now been completed and extended to ten paired seeds.
heterogeneity makes $B_r$ strongly more adverse at both values of $E$.
Increasing $E$ changes $L_r$ in a regime-dependent direction: it becomes more
negative under IID data but strongly positive under strong non-IID data.  In
the strong non-IID, $E=5$ cell, first-order alignment is positive in $99.0\%$
of audited snapshots while the objective decreases in only $69.7\%$ of them.
This confirms that a useful analysis must retain both cancellation among
$P_r,L_r,B_r$ and the finite-step curvature contribution.  All key directional
contrasts hold for all ten seed pairs, so the mechanism audit is sufficiently
stable to motivate a formal covariance assumption.
